\section{Experiments}
\label{experiments}

\subsection{Research Questions}
\label{sec:exp-rq}

The evaluation is organized around the following research questions:
\begin{description}[leftmargin=*,style=nextline]
  \item[RQ1] How accurately can the approach transform natural language requirements into executable tasks?
  \item[RQ2] Does knowledge guidance improve alignment and execution completeness?
  \item[RQ3] Does multi-agent task planning improve planning consistency compared with simpler transformation baselines?
  \item[RQ4] How does human review affect the quality of the final executable requirement representation?
\end{description}

\subsection{Experimental Setup}
\label{sec:exp-setup}

The experiment evaluates the complete requirement-to-task workflow.
For each requirement, the approach produces semantic analysis artifacts, task decompositions, action sequences, requirement-to-task mappings, and execution flows.
All outputs are compared against manually inspected reference artifacts or expert judgments.

\subsection{Datasets}
\label{sec:exp-datasets}

The dataset should contain natural language requirements from course project specifications or representative intelligent software system scenarios.
Each requirement should be annotated with expected task units, action semantics, and dependency information.
This section will report dataset sources, selection criteria, annotation protocol, and basic statistics once the dataset is finalized.

\subsection{Baselines}
\label{sec:exp-baselines}

The planned baselines are:
\begin{itemize}[leftmargin=*]
  \item Direct prompting that asks an LLM to produce tasks from a requirement in a single step.
  \item A single-agent planner without role separation.
  \item A template-only transformation method without knowledge-guided reasoning.
  \item A variant without human review.
\end{itemize}

\subsection{Evaluation Metrics}
\label{sec:exp-metrics}

\subsubsection{Task Generation Accuracy}
\label{sec:metric-accuracy}

Task generation accuracy measures whether generated tasks correspond to valid actions implied by the original requirement.

\subsubsection{Execution Completeness}
\label{sec:metric-completeness}

Execution completeness measures whether the task workflow covers the necessary steps, inputs, outputs, and dependencies required to operationalize the requirement.

\subsubsection{Requirement-to-Task Alignment}
\label{sec:metric-alignment}

Requirement-to-task alignment measures whether each generated task can be traced to a specific requirement fragment without introducing unsupported behavior.

\subsubsection{Plan Consistency}
\label{sec:metric-consistency}

Plan consistency measures whether the action sequence respects dependency order, constraints, and expected flow.

\subsection{Main Results}
\label{sec:exp-results}

This subsection will report quantitative and qualitative results for the research questions.
No synthetic performance numbers are included in this structure draft.

\subsection{Ablation Study}
\label{sec:exp-ablation}

The ablation study will evaluate the contribution of knowledge guidance, agent role separation, planning rules, task templates, and human review.
Each ablation removes one component while keeping the remaining workflow unchanged.

\subsection{Case Study}
\label{sec:exp-case}

The case study will show one representative requirement and trace its transformation into semantic elements, tasks, action sequences, requirement-to-task mappings, and an execution flow.
